% !TEX root = ../Report.tex
% ======================================================================
\chapter{Introduction}
\label{ch:intro}
% ======================================================================

This report presents the design, implementation, testing, and evaluation of GiTrip, a web application for collaborative trip planning via Git~\cite{ref_git}-style version control system (VCS) combined with constraints-aware auto-scheduler.

\section{Problem Context and Motivation}
Trip planning is time-consuming because a feasible schedule must satisfy
multiple constraints simultaneously, including travel time between destinations,
opening hours, stay durations, and daily activity limits. The process becomes
more complex for group trips, where participants propose alternative routes and
schedules, with edits often distributed across chat messages, screenshots, and
shared documents. This fragmented approach increases the
likelihood of duplicated work, accidental overwrites, and loss of rationale for
changes.

While existing route planners can optimise a single itinerary and generic collaboration tools can support shared documents, optimising the entire trip schedule and supporting change management in group trip planning remain insufficient.

To validate the practical relevance of this problem, a preliminary survey was
conducted with 30 participants (full questionnaire and results in Appendix~\ref{appx:prelim_survey}).
Key findings support three observations:
\begin{enumerate}
  \item \textbf{Planning is widely effortful:} 60\% (18/30) reported trip
        planning as \textit{often} or \textit{always} stressful/time-consuming; 93\% (28/30) as
        \textit{at least sometimes} stressful.
  \item \textbf{Fragmented tool usage is common:} 83\% (25/30) use Google
        Maps~\cite{ref_google_maps} for planning; 47\% (14/30) use Google Docs~\cite{ref_google_docs}/Microsoft Word~\cite{ref_ms_word}
        (multi-select). Current practice fragments planning across maps and
        general-purpose documents.
  \item \textbf{Demand for an integrated solution:} 90\% (27/30) responded
        \textit{very likely} or \textit{somewhat likely} to use GiTrip.
\end{enumerate}
These findings, initially presented in the Progress Report, motivated the
background research and competitor analysis in
Chapter~\ref{ch:background}.

\section{Project Overview}
GiTrip applies Git-style VCS semantics to trip planning. A trip is represented as a repository containing the users’ commits (changes made to the trip). Users can branch (make alternative routes) and merge branches via structured conflict resolution to make the final trip plan. This version-controlled plan is integrated with an automatic scheduler that produces a non-overlapping itinerary while respecting constraints such as active hours, per-stop stay duration, user preferences, and opening hours.

The key design goal is usability for non-experts. Git concepts are surfaced through a clear web interface (buttons for branching/merging, workflow graph, conflict resolver screens). Trip plans are visualised using an interactive timeline (editable by dragging) and a map view showing ordered stops and route geometries. The ``Git mode'' toggle, when turned off, hides Git-specific vocabulary for non-technical users. The feature-ranking results (Q5) of the survey reinforce this requirement: concrete itinerary features such as timeline/map visualisation, travel-time estimation, and automatic optimisation are perceived as most helpful, while Git-style trip repository concepts are less directly valued by some respondents (Appendix~\ref{app:agg-res}). Therefore, Git-style semantics should be implemented to enable VCS, but presented through an approachable interface that does not assume prior Git knowledge.

\section{Objectives and Significance}

The objectives, categorised by MoSCoW (Must/Should/Could/Won't) prioritisation, were initially formalised in the Project Specification. Must-have objectives include covering repository history,
an optimised automatic scheduler, routing integration, interactive
visualisation, and merging with conflict resolution. The full requirements and
their rationale are detailed in Chapter~\ref{ch:requirements}.

GiTrip is a suitably challenging 3rd-year Computer Science project because it combines: (i) VCS semantics over structured data, (ii) constraint-aware scheduling with external routing services, and (iii) interactive visualisation with editable timelines and maps. The project makes a novel contribution by applying VCS semantics to collaborative trip planning, addressing a documented gap in existing tools that are further discussed in Chapter~\ref{ch:background}. This is also significant because the combination of constraint-aware scheduling with structured change management has potential applications beyond tourism when generalised, including project planning and resource allocation scenarios.